\documentclass[UTF8, a4paper, 11pt]{ctexart}

% --- 核心宏包 ---
\usepackage[top=2.38cm, bottom=2cm, left=2cm, right=2cm]{geometry}
\usepackage{tabularx}
\usepackage{amssymb}
\usepackage{setspace}
\usepackage{enumitem}
\usepackage{tcolorbox}
\tcbuselibrary{skins, breakable}
\usepackage{multirow}

% --- 设置行间距（1.2倍） ---
\setstretch{1.2}

% --- 定义统一的双线边框环境（距上下2cm，可跨页） ---
\newtcolorbox{UnifiedBox}[1][]{
    enhanced,
    breakable,
    sharp corners,
    colback=white,
    colframe=black,
    boxrule=0.6pt,
    borderline={0.6pt}{-3pt}{black},
    left=0.5em, right=0.5em, top=0.5em, bottom=0.5em,
    width=\textwidth,
    height fill=maximum,
    pad at break=3mm,
    toprule at break=0.6pt,
    bottomrule at break=0.6pt,
    enlargepage flexible=\baselineskip,
    skin first=enhanced,
    skin middle=enhanced,
    skin last=enhanced,
    extras first={
        borderline south={0.6pt}{-3pt}{black},
    },
    extras middle={
        borderline north={0.6pt}{-3pt}{black},
        borderline south={0.6pt}{-3pt}{black},
    },
    extras last={
        borderline north={0.6pt}{-3pt}{black},
        borderline south={0.6pt}{-3pt}{black},
    },
    underlay={
        \node[anchor=south, yshift=-0.6cm] at (frame.south) {\zihao{5} \thepage};
    },
    #1
}

\begin{document}

% ===================== 第 1 页：封面 =====================
\begin{center}
    {\zihao{2} \bfseries 专业学位硕士研究生 \\ \medskip 毕业（学位）论文开题报告}
\end{center}

\vspace{0.5cm}

\noindent
\begin{tcolorbox}[
    enhanced,
    sharp corners,
    colback=white,
    colframe=black,
    boxrule=0.6pt,
    left=0pt, right=0pt, top=0pt, bottom=0pt,
    boxsep=0pt,
    outer arc=0pt,
    arc=0pt
]
\resizebox{\linewidth}{!}{%
\begin{tabular}{|>{\centering\arraybackslash}p{1.6cm}|>{\centering\arraybackslash}p{2.8cm}|>{\centering\arraybackslash}p{1.6cm}|>{\centering\arraybackslash}p{2.8cm}|>{\centering\arraybackslash}p{1.6cm}|>{\centering\arraybackslash}p{2.8cm}|}
\hline
\bfseries 姓\quad 名 & & \bfseries 学\quad 号 & & \bfseries 已修学分 & \\ \hline
\bfseries 所属学院 & \multicolumn{2}{c|}{工商管理学院} & \multicolumn{2}{c|}{\bfseries 学科专业} & 工商管理硕士（MBA） \\ \hline
\bfseries 指导教师 & \multicolumn{2}{c|}{} & \multicolumn{2}{c|}{\bfseries 选题时间} & 2025 年 \quad 月 \quad 日 \\ \hline
\bfseries 研究方向 & \multicolumn{5}{l|}{
    \begin{minipage}[c][2.2cm]{0.72\linewidth}
    \fontsize{7.5pt}{9pt}\selectfont \vspace{5pt}
    $\square$运营与项目管理 \quad $\square$组织行为与人力资源管理 \quad $\square$市场营销与电子商务 \quad $\square$公司理财 \\ \medskip
    $\square$金融管理与金融科技 \quad $\square$企业战略与创新创业管理 \quad $\square$大数据与信息管理 \quad $\square$国际商务管理 \vspace{5pt}
    \end{minipage}
} \\ \hline
\bfseries 论文题目 & \multicolumn{5}{c|}{华润数字科技有限公司发展战略研究} \\ \hline
\end{tabular}%
}
\end{tcolorbox}

\vspace{2em}
\noindent {\zihao{4} \bfseries 一、选题背景及意义}

\vfill
\begin{center} 1 \end{center}
\newpage

% ===================== 第 2-3 页：课题来源、选题的目的与意义 =====================
\clearpage
\setcounter{page}{2}

\begin{UnifiedBox}
\zihao{-4}
\setstretch{1.6}
\setlength{\parindent}{2em}

\noindent \textbf{（课题来源、选题的目的与意义）}

\vspace{0.2cm}
\noindent \textbf{（一）课题来源}

\vspace{0.2cm}
当前，我们正处在一个由数字技术驱动的新一轮科技革命与产业变革的历史交汇期。大数据、云计算、人工智能、区块链等新一代信息技术正以前所未有的力量重塑全球经济格局与产业生态。在此背景下，数字化转型已不再是企业的可选项，而是关乎生存与发展的必由之路。我国高度重视数字经济发展，在国家“十四五”规划纲要中明确提出“加快数字化发展，建设数字中国”的战略部署，推动数字经济和实体经济深度融合，打造具有国际竞争力的数字产业集群。这一系列国家战略为数字科技企业的发展提供了广阔的市场空间和历史性的机遇。

华润数字科技有限公司作为大型央企——华润集团旗下专注数字科技业务的旗舰公司，肩负着推动集团整体数字化转型、打造新型数字经济体、并对外赋能产业数字化重要使命。公司业务涵盖云计算、大数据、人工智能、物联网等多个前沿领域，是华润集团拥抱数字时代、培育新质生产力的关键载体。从本质上而言，数字科技行业属于技术密集、创新驱动型行业，其发展对技术前瞻性、战略敏捷性及商业模式创新性的要求极高，其对核心竞争力的依赖强于传统行业。因此，制定一套科学、清晰、前瞻的发展战略对于华润数科而言显得至关重要。一套合适的发展战略能够指引公司在复杂多变的市场环境中明确方向、配置资源、构筑壁垒，是实现其战略愿景的根本保障。

在实践过程中，华润数科在快速发展的同时也面临着一系列严峻的挑战。一方面，外部技术迭代迅速，市场竞争日趋激烈，不仅面临其他科技巨头的挤压，还需应对众多灵活创新的科技初创公司的挑战；另一方面，公司自身在战略聚焦、核心能力构建、创新机制与管理模式等方面仍存在优化空间。公司管理层虽然已经确立了初步的战略方向，但在具体落地路径、业务协同、以及如何将集团资源优势转化为市场胜势等方面，仍需进一步探索与细化。基于此背景下，本文将围绕企业战略管理的核心维度，对华润数科的内外部环境进行系统剖析，旨在为其制定一个既能契合国家战略导向、又能充分发挥自身优势的可持续发展战略。

\vspace{0.3cm}
\noindent \textbf{（二）选题目的}

\vspace{0.2cm}
企业的持续健康发展与科学战略的引领密不可分。本文以华润数科为研究对象，在研究过程中旨在系统分析其发展战略现状及内外部环境，指出当前制约公司高质量发展的关键问题，并为华润数科未来的战略制定与实施提供具体可行的参考方案。选题目的主要有以下几点：

（1）通过发展战略的深入研究，可以帮助华润数科明确其在数字经济浪潮中的核心定位与主攻方向，使公司的技术研发、产品布局、市场拓展等活动能够统一在清晰的战略目标之下，满足公司自身跨越式发展的内在需要。通过战略的明晰，公司上下能够统一思想、凝聚共识，更深刻地理解其肩负的使命与愿景。

（2）通过科学的发展战略研究，可以为华润数科识别并培育新的业务增长点，优化现有业务结构，建立起持续创新的长效机制。这能够帮助公司在激烈的市场竞争中构筑坚实的护城河，应对未来不确定性，实现从“跟随者”向“引领者”的转变，满足华润集团对其作为新增长引擎的战略期待。

（3）从更宏观的视角来看，一个成功的华润数科发展战略，不仅能够驱动华润集团整体的数字化转型与产业升级，更能作为央企发展数字科技的典范，在落实国家数字中国战略、赋能实体经济数字化转型过程中，充分发挥国有经济的主导作用和战略支撑力量，实现企业价值与社会价值的统一。

\vspace{0.4cm}
\noindent \textbf{（三）选题意义}

\vspace{0.2cm}
制定科学、系统、前瞻的发展战略，目的是要为企业指明在复杂环境下的前进路径，有效配置有限资源，从而在激烈的市场竞争中脱颖而出，实现可持续的价值创造。而通过系统性的战略分析与规划，属于企业为了实现长期目标而采取的核心管理活动，其目的是提升企业的整体效能，使企业的资源能力与环境机会相匹配，是为了让企业更适应未来的挑战，能够把握时代机遇，也是为了满足国家战略与产业升级对领军企业的呼唤，达到更高的组织效能，促进企业宏伟蓝图的实现。企业发展战略研究在企业管理学科中占有举足轻重的地位，对企业的长远发展有着非凡的意义。本文的研究意义有以下几方面：

（1）实际意义。论文以华润数科作为实际案例进行研究，系统梳理并诊断其当前发展战略中存在的优势与不足，结合公司面临的宏观环境与行业趋势，提出适合华润数科的、可落地的战略优化方案与实施路径，为公司的战略决策与执行提供直接参考，具有显著的实际应用价值。

（2）理论意义。本文将综合运用战略管理的经典理论，并结合数字科技企业的新特性，在真实的商业场景中进行验证与拓展。通过理论联系实际，探索数字经济时代下企业战略管理的新模式与新特征，对丰富和发展战略管理理论在数字产业领域的应用具有一定的理论意义。

（3）借鉴意义。论文的研究成果可为其他正处于数字化转型中的大型企业集团，特别是央企和国企，在如何孵化、培育和壮大数字科技业务板块方面提供宝贵的借鉴。在当今百年未有之大变局下，传统企业如何利用数字科技实现“第二曲线”增长，如何应对科技公司的跨界竞争，论文所提出的战略框架与思考能为其提供有益的启示和参考。

\vfill
\begin{flushright}
    （本表可附页）
\end{flushright}
\end{UnifiedBox}
\newpage

% ===================== 第 4-7 页：文献综述 =====================
\vspace{0.6cm}
\noindent {\zihao{4} \bfseries 二、文献综述}

\begin{UnifiedBox}
\zihao{-4}
\setstretch{1.6}
\setlength{\parindent}{2em}

\noindent \textbf{（与选题有关的国内外发展现状和动态）}

\vspace{0.3cm}
\noindent \textbf{（一）公司发展战略理论演进与核心框架研究}

\vspace{0.2cm}
公司发展战略理论历经数十年迭代，已形成涵盖战略制定、实施与风险管理的完整体系，为企业战略实践提供了坚实理论支撑。杨华江（2021）通过梳理公司战略和战略风险管理理论的演进脉络，指出战略理论已从早期``资源基础观''''产业结构观''，逐步向``动态能力理论''''生态系统理论''转型，尤其在数字经济背景下，企业战略更强调对外部环境变化的快速响应与内部资源的动态整合。这一演进趋势为华润数字科技有限公司这类科技型企业的战略制定提供了理论坐标系，其战略设计需兼顾技术迭代速度与市场需求变化的双重不确定性。

在战略分析工具应用方面，现有研究形成了以SWOT分析为基础、多工具协同的方法论体系。李亚祥（2024）以铁路物资公司为研究对象，通过SWOT分析法系统拆解企业内部优势（如供应链整合能力）、劣势（如数字化水平不足）与外部机会（如政策支持）、威胁（如行业竞争加剧），最终提出``数字化转型+差异化服务''的战略组合。汪恭二、占寿罡、万超（2023）进一步优化SWOT分析框架，引入权重赋值法量化各因素影响程度，为T公司制定了优先级明确的战略实施路径。此外，孙虹（2020）提出的IFE-EFE矩阵分析法，通过量化评估企业内部因素（如技术研发能力）与外部因素（如市场增长率）的权重得分，为X公司构建了``战略匹配度-实施可行性''双维度评估模型，该方法对华润数科这类需平衡技术投入与市场回报的企业具有直接借鉴意义。

五力模型作为行业竞争环境分析的核心工具，在跨行业战略研究中得到广泛应用。袁琪（2021）基于五力模型分析大连机场货运公司的竞争格局，指出其需通过``提升供应商议价能力（如与航空公司签订长期协议）+降低潜在进入者威胁（如构建一体化物流网络）''的组合策略应对行业竞争。虽然该研究聚焦传统服务业，但其中``通过资源整合构建竞争壁垒''的思路，可为华润数科在数字服务领域（如政企数字化转型服务）的战略布局提供参考，即通过整合华润集团产业资源，形成差异化竞争优势。

\vspace{0.3cm}
\noindent \textbf{（二）跨行业公司发展战略实践研究}

\vspace{0.2cm}
\noindent \textbf{1.传统行业数字化转型战略研究}

传统行业企业的数字化转型战略研究，为华润数科探索``产业数字化服务''路径提供了实践参考。祝志华（2024）以T汽车零部件公司为案例，提出传统制造企业需以``技术赋能+业务重构''为核心，通过引入工业互联网平台、优化供应链数字化管理，实现从``生产导向''向``服务导向''的战略转型。王广昀、王凤兰、赵波等（2021）针对大庆油田公司的研究则进一步指出，资源型企业的数字化战略需聚焦``核心业务场景穿透''，如通过勘探开发数字化平台提升资源开采效率，同时构建``数字化+绿色化''双轮驱动模式，这一思路与华润数科服务能源、制造等传统产业的业务定位高度契合。

Shuai W、Hongxiang C（2017）以Wugang公司为核心的平顶山钢铁产业集群研究，提出``龙头企业引领+产业链协同''的数字化转型模式，即由龙头企业搭建产业互联网平台，带动上下游中小企业实现数字化升级。该模式对华润数科而言具有重要借鉴价值，其可依托华润集团在零售、能源、医药等领域的产业优势，构建``产业数字化平台+垂直场景解决方案''的服务模式，推动集团内部产业协同与外部行业赋能的双向发展。

\vspace{0.2cm}
\noindent \textbf{2.科技与互联网行业创新战略研究}

科技与互联网行业的创新战略研究，为华润数科的技术研发与商业模式设计提供了直接参考。龙敏艳（2022）对华为的案例研究表明，科技型企业需以``持续研发投入+全球化布局''为战略核心，通过构建``基础研究-应用开发-市场转化''的全链条创新体系，实现技术领先与市场份额的双重突破。虽然华为聚焦ToC端消费电子与通信设备领域，但其``以技术创新驱动战略迭代''的逻辑，可迁移至华润数科ToB端数字服务业务，如通过加大政企数字化转型解决方案的研发投入，形成技术壁垒。

Vochozka H（2021）从可持续发展视角指出，科技企业的创新战略需兼顾``技术突破''与``社会责任''，通过绿色技术研发、数据安全保障等举措，实现商业价值与社会价值的统一。这一观点对华润数科具有特殊意义，其作为国有科技企业，需在战略设计中融入``国有资本使命''，如在数据安全、数字乡村建设等领域主动承担社会责任，形成``商业价值-社会价值''协同共生的战略格局。

Clayton AP（2018）以美国北卡罗来纳州无线研究中心为案例，提出``政产学研协同创新''的战略模式，即通过政府政策引导、企业资金支持、高校科研攻关的三方协同，加速技术成果转化。该模式可为华润数科探索``技术生态构建''提供思路，其可联合高校、科研机构共建数字技术实验室，聚焦工业互联网、人工智能等核心领域，推动技术研发与产业应用的深度融合。

\vspace{0.2cm}
\noindent \textbf{3.服务行业平台化战略研究}

服务行业的平台化战略研究，为华润数科构建数字服务生态提供了实践经验。陈凯麟、蒋伏心（2017）基于``供给侧改革''背景，对GL证券公司的研究表明，金融服务企业需以``互联网平台+场景嵌入''为核心，通过搭建一站式金融服务平台、嵌入客户日常消费场景，实现服务模式创新。刘建刚、钱玺娇（2016）对小米科技的案例研究进一步提出``技术创新+商业模式创新''双轮驱动模型，指出小米通过``硬件引流+互联网服务变现''的模式，构建了多元化盈利体系。虽然上述研究聚焦金融与消费电子领域，但其``平台化整合+场景化服务''的思路，可迁移至华润数科的数字服务业务，如搭建政企数字化服务平台，将数据服务、云服务等嵌入客户业务场景，实现从``单一服务提供商''向``生态型服务平台''的战略转型。

\vspace{0.3cm}
\noindent \textbf{（三）现有研究不足与本文研究方向}

\vspace{0.2cm}
综合现有研究来看，公司发展战略研究已形成``理论-工具-实践''三位一体的体系，但针对国有背景科技企业的战略研究仍存在以下不足：

（1）研究对象聚焦偏差。现有研究多集中于传统制造企业（如祝志华，2024）、消费端科技企业（如龙敏艳，2022）或金融服务企业（如陈凯麟、蒋伏心，2017），对兼具``国有资本属性''''产业数字化服务定位''''集团协同发展需求''的企业关注不足。华润数科作为华润集团数字化转型的核心载体，其战略制定需同时满足``国有资产保值增值''''集团产业赋能''''市场化竞争''三重目标，现有研究难以提供针对性的理论与实践支撑。

（2）研究视角存在局限。现有研究多从单一维度（如技术创新、市场拓展）展开战略分析，缺乏对``集团资源协同-行业需求匹配-技术能力支撑''多维度的系统整合。例如，Vochozka H（2021）虽强调科技创新与可持续发展的协同，但未涉及集团层面的资源整合机制；Shuai W、Hongxiang C（2017）关注产业集群协同，却未深入探讨科技企业如何平衡集团内部资源调用与外部市场竞争。华润数科的战略设计需兼顾``调用集团零售、能源等产业资源''''满足政企客户数字化需求''''提升自身技术研发能力''的三重协同，现有研究的单一视角难以覆盖这一复杂需求。

（3）研究方法有待深化。现有研究多采用``案例分析+定性描述''的方法，缺乏对``战略实施效果-影响因素''的量化分析。例如，李亚祥（2024）、汪恭二、占寿罡、万超（2023）的研究虽通过SWOT分析提出战略方向，但未建立战略实施效果的评估指标体系；孙虹（2020）的IFE-EFE矩阵分析法虽实现部分量化，但未结合企业具体业务场景设计动态调整机制。华润数科的战略研究需构建``量化评估-动态调整''的闭环体系，以应对数字技术快速迭代与市场需求变化的挑战。

基于上述不足，本文以华润数科为研究对象，将从以下三方面展开研究：

多维度战略环境分析：结合PEST分析法与SWOT-权重赋值模型，系统拆解华润数科面临的政策环境、技术环境、市场环境，以及内部资源能力，形成``内外部因素-影响程度-战略匹配度''的三维分析框架。

集团协同型战略设计：基于华润集团``多元化产业+国有资本''的独特优势，设计``集团资源整合-行业场景赋能-技术创新驱动''的协同战略，具体包括``依托集团产业场景培育数字服务能力''''通过技术输出反哺集团数字化转型''''联合外部生态伙伴拓展市场空间''三大路径。

动态化战略评估体系：引入平衡计分卡（BSC）与KPI结合的评估模型，从``财务指标、客户指标、内部流程指标、学习成长指标''四个维度，构建战略实施效果的动态评估机制，确保战略能够根据内外部环境变化及时调整。

\vfill
\begin{flushright} （本表可附页） \end{flushright}
\end{UnifiedBox}

% ===================== 第 7 页下方：三、研究的主要内容及技术路线 =====================
\vspace{0.6cm}
\noindent {\zihao{4} \bfseries 三、研究的主要内容及技术路线}

\newpage

% ===================== 第 8-13 页：合并为一个连续方框 =====================
\begin{UnifiedBox}
\zihao{-4}
\setstretch{1.6}
\setlength{\parindent}{2em}

\noindent \textbf{（一）研究内容及核心方向}

\vspace{0.2cm}
\noindent \textbf{1.研究内容}

\vspace{0.2cm}
\noindent \textbf{（1）华润数科内外部战略环境定性分析}

\noindent \textbf{外部环境分析：}结合PEST分析法与五力模型，从政策、技术、市场三个核心维度展开。政策层面，梳理``数字中国''''国企数字化转型''等国家政策及粤港澳大湾区、长三角等地方数字产业扶持政策，明确政策导向带来的发展机遇与合规约束；技术层面，参考Gartner、IDC等权威机构技术报告，分析人工智能、工业互联网、数据安全等核心技术的应用成熟度与场景边界，判断技术迭代对业务布局的影响；市场层面，整合艾瑞咨询、头豹研究院等行业报告，梳理政企数字化转型市场的细分需求与竞争格局，通过头部企业公开战略与产品信息，总结其竞争优势与短板，明确华润数科的市场定位空间。

\noindent \textbf{内部资源能力梳理：}从资源与能力两大维度搭建分析框架。资源维度，依托华润集团年度报告、华润数科公开信息及专利局公示数据，梳理可复用的集团资源、自有技术资产与财务基础；能力维度，通过公开项目案例、客户评价及合作伙伴信息，分析技术研发、项目交付、生态协同三大能力现状，总结``集团产业场景积淀''等核心优势与``市场化激励机制不足''等突出短板。

\vspace{0.3cm}
\noindent \textbf{（2）华润数科集团协同型发展战略设计}

\vspace{0.2cm}
\noindent \textbf{战略定位与目标体系：}结合内外部环境分析结果，参考中国电子科技集团等国有科技企业数字业务的战略定位经验，确定``集团数字化转型核心载体+行业数字服务生态构建者''的双重定位。目标体系分阶段设定：短期（1-2年）聚焦集团内部数字化服务覆盖，优先完成华润零售、能源板块的核心业务数字化改造；中期（3-5年）拓展外部行业客户，形成``集团业务+外部市场''双轮驱动格局；长期（5年以上）打造国内产业数字服务领域标杆，形成差异化竞争优势。

\noindent \textbf{核心战略路径规划：}围绕``集团资源协同''''技术创新''''生态构建''三大方向设计具体路径。集团资源协同路径，依据华润集团内部资源共享制度，明确零售场景数据、能源业务经验等资源的调用流程与收益分配规则，实现集团内资源的高效复用；技术创新路径，结合行业技术发展趋势与自身业务需求，确定工业互联网平台、数据安全解决方案两大核心技术方向，分阶段推进技术研发与产品落地；生态构建路径，参考树根互联、用友等企业的生态合作模式，筛选高校科研机构、第三方技术公司、行业协会等合作伙伴，建立``联合研发+资源互补''的合作机制，拓展服务边界。

\noindent \textbf{（3）战略实施保障体系构建}

\noindent \textbf{合规与激励保障：}针对国有科技企业特性，建立双重保障机制。合规层面，依据《数据安全法》《国有企业合规管理办法》及集团内控要求，制定数据采集、存储、应用全流程的合规审查流程，确保业务开展符合政策与监管要求；激励层面，参考科创板国企股权激励、项目跟投等公开案例，在``创新业务板块''试点灵活激励机制，同时依据国资委容错政策，明确创新失败的认定标准与免责范围，激发团队创新活力。

\noindent \textbf{资源与评估保障：}资源保障方面，结合公司近3年研发、人才、资金投入情况，制定未来三年资源规划，优先保障核心技术研发与关键人才引进；评估保障方面，基于平衡计分卡思路，从财务、客户、内部流程、学习成长四个维度设定评估指标，通过定期复盘跟踪战略落地效果，及时调整执行策略。

\vspace{0.3cm}
\noindent \textbf{2.核心研究方向}

聚焦``国有集团背景下数字科技企业的资源协同与市场化发展''这一核心命题，重点解决三个关键问题：一是如何高效复用集团产业资源，形成差异化服务能力；二是如何平衡国有资产监管要求与市场化经营灵活性；三是如何在竞争激烈的数字服务市场中找到可持续的业务增长点。

\vspace{0.5cm}
\noindent \textbf{（二）拟采取的研究方法、技术路线及可行性分析}

\vspace{0.2cm}
\noindent \textbf{1.拟采取的研究方法}

\vspace{0.2cm}
\noindent \textbf{（1）文献研究法}

系统梳理公司发展战略理论、国有科技企业转型案例及数字服务行业报告，参考杨华江（2021）关于战略理论演进的研究、刘建刚（2016）平台化战略案例分析等成果，为研究提供理论支撑与经验借鉴，确保战略设计的理论严谨性。

\vspace{0.2cm}
\noindent \textbf{（2）案例分析法}

选取三类标杆案例深入研究：一是国有科技企业，通过其公开战略规划、年度报告，总结国有背景下数字业务的资源整合与市场化运作经验；二是产业数字服务厂商，依据产品白皮书、合作案例，分析技术路线选择与生态构建模式；三是华润集团内部数字化案例，通过集团官网披露信息、项目新闻报道，提炼集团资源复用的可行路径，为战略设计提供直接参考。

\vspace{0.2cm}
\noindent \textbf{（3）定性分析法}

通过逻辑推导与经验总结，对华润数科内外部环境、战略路径进行定性判断。例如，基于政策文本解读明确政策对业务的影响方向，结合市场竞争格局判断差异化定位的可行性，参考同类企业经验设计战略实施保障机制，确保研究结论贴合企业实际经营场景。

\vspace{0.3cm}
\noindent \textbf{2.技术路线}

\vspace{0.2cm}
\noindent \textbf{（1）研究启动阶段（第1-2个月）}

明确研究范围：确定以华润数科核心业务为研究对象，排除非核心业务。

资料收集：通过华润集团官网、行业数据库、学术期刊收集政策文件、行业报告、理论文献与案例资料，完成文献综述与案例库搭建。

\vspace{0.2cm}
\noindent \textbf{（2）环境分析阶段（第3-4个月）}

外部环境研究：通过政策文本解读、技术报告分析、市场信息整理，形成PEST分析结论；结合竞争对手公开信息，完成五力模型竞争格局梳理，明确外部机遇与威胁。

内部资源能力梳理：基于华润数科公开信息与内部资料，整理资源与能力现状，总结优势与短板，形成内部分析报告。

\vspace{0.2cm}
\noindent \textbf{（3）战略设计阶段（第5-6个月）}

战略定位研讨：结合内外部分析结果与标杆案例经验，确定双重战略定位与分阶段目标。

路径与保障设计：参考同类企业经验，规划集团资源协同、技术创新、生态构建三大路径，同步设计合规、激励、资源保障机制，形成战略方案初稿。

\vspace{0.2cm}
\noindent \textbf{（4）验证与完善阶段（第7-8个月）}

战略验证：选取华润数科现有业务，结合项目运行实际情况，验证战略路径的可行性。

方案修订：根据验证结果，调整战略目标与实施路径，形成最终战略方案与研究报告。

\vspace{0.3cm}
\noindent \textbf{3.可行性分析}

\vspace{0.2cm}
\noindent \textbf{（1）理论可行性}

现有公司发展战略理论已成熟，国有科技企业转型、数字服务行业发展等领域的研究成果可直接借鉴，为战略设计提供充足理论支撑；同时，PEST、五力模型等分析工具在定性研究中应用广泛，方法体系成熟可靠。

\vspace{0.2cm}
\noindent \textbf{（2）资料可行性}

外部资料方面，政策文件、行业报告、竞争对手公开信息可通过政府官网、商业数据库、权威媒体等合法渠道获取，信息覆盖全面；内部资料方面，华润集团与华润数科的公开报告、专利信息、项目案例等可直接用于内部资源能力分析，若需进一步补充细节，可通过合规沟通获取，资料获取难度较低。

\vspace{0.5cm}
\noindent \textbf{（三）论文框架}

\vspace{0.2cm}
\noindent \textbf{第一章\quad 绪论}

\hangindent=2em \hangafter=1
1.1 研究背景与意义

\hangindent=4em \hangafter=1
\quad 1.1.1 研究背景

\hangindent=4em \hangafter=1
\quad 1.1.2 研究意义

\hangindent=2em \hangafter=1
1.2 国内外研究现状

\hangindent=4em \hangafter=1
\quad 1.2.1 国外研究现状

\hangindent=4em \hangafter=1
\quad 1.2.2 国内研究现状

\hangindent=4em \hangafter=1
\quad 1.2.3 研究评述

\hangindent=2em \hangafter=1
1.3 研究内容、方法与技术路线

\hangindent=4em \hangafter=1
\quad 1.3.1 研究内容

\hangindent=4em \hangafter=1
\quad 1.3.2 研究方法

\hangindent=4em \hangafter=1
\quad 1.3.3 技术路线

\hangindent=2em \hangafter=1
1.4 研究创新点

\vspace{0.3cm}
\noindent \textbf{第二章\quad 相关理论基础与分析工具}

\hangindent=2em \hangafter=1
2.1 核心理论支撑

\hangindent=4em \hangafter=1
\quad 2.1.1 动态能力理论

\hangindent=4em \hangafter=1
\quad 2.1.2 生态系统理论

\hangindent=4em \hangafter=1
\quad 2.1.3 国有资产管理理论

\hangindent=2em \hangafter=1
2.2 战略分析工具

\hangindent=4em \hangafter=1
\quad 2.2.1 PEST分析法

\hangindent=4em \hangafter=1
\quad 2.2.2 五力模型

\hangindent=4em \hangafter=1
\quad 2.2.3 SWOT分析法

\hangindent=2em \hangafter=1
2.3 数字服务行业特性

\hangindent=4em \hangafter=1
\quad 2.3.1 技术驱动与场景依赖的双重属性

\hangindent=4em \hangafter=1
\quad 2.3.2 国有数字科技企业的战略特殊性

\vspace{0.3cm}
\noindent \textbf{第三章\quad 华润数科内外部战略环境分析}

\hangindent=2em \hangafter=1
3.1 外部环境分析

\hangindent=4em \hangafter=1
\quad 3.1.1 政策环境

\hangindent=4em \hangafter=1
\quad 3.1.2 技术环境

\hangindent=4em \hangafter=1
\quad 3.1.3 市场环境

\hangindent=2em \hangafter=1
3.2 内部资源能力分析

\hangindent=4em \hangafter=1
\quad 3.2.1 资源维度

\hangindent=4em \hangafter=1
\quad 3.2.2 能力维度

\hangindent=4em \hangafter=1
\quad 3.2.3 优势与短板总结

\hangindent=2em \hangafter=1
3.3 内外部环境匹配（SWOT分析）

\hangindent=4em \hangafter=1
\quad 3.3.1 优势-机会（SO）策略方向

\hangindent=4em \hangafter=1
\quad 3.3.2 优势-威胁（ST）应对思路

\hangindent=4em \hangafter=1
\quad 3.3.3 劣势-机会（WO）改进路径

\hangindent=4em \hangafter=1
\quad 3.3.4 劣势-威胁（WT）风险规避

\vspace{0.3cm}
\noindent \textbf{第四章\quad 华润数科集团协同型发展战略设计}

\hangindent=2em \hangafter=1
4.1 战略定位与目标

\hangindent=4em \hangafter=1
\quad 4.1.1 双重定位

\hangindent=4em \hangafter=1
\quad 4.1.2 分阶段目标

\hangindent=2em \hangafter=1
4.2 核心战略路径

\hangindent=4em \hangafter=1
\quad 4.2.1 集团资源协同路径

\hangindent=4em \hangafter=1
\quad 4.2.2 技术创新路径

\hangindent=4em \hangafter=1
\quad 4.2.3 生态构建路径

\hangindent=2em \hangafter=1
4.3 细分业务策略

\hangindent=4em \hangafter=1
\quad 4.3.1 集团内部数字化服务

\hangindent=4em \hangafter=1
\quad 4.3.2 外部行业数字服务

\vspace{0.3cm}
\noindent \textbf{第五章\quad 华润数科战略实施保障体系}

\hangindent=2em \hangafter=1
5.1 合规保障机制

\hangindent=4em \hangafter=1
\quad 5.1.1 数据安全管理流程

\hangindent=4em \hangafter=1
\quad 5.1.2 国资监管对接

\hangindent=2em \hangafter=1
5.2 激励与资源保障

\hangindent=4em \hangafter=1
\quad 5.2.1 市场化激励机制

\hangindent=4em \hangafter=1
\quad 5.2.2 资源投入规划

\hangindent=2em \hangafter=1
5.3 战略评估与调整

\hangindent=4em \hangafter=1
\quad 5.3.1 评估指标设计

\hangindent=4em \hangafter=1
\quad 5.3.2 评估与调整流程

\vspace{0.3cm}
\noindent \textbf{第六章\quad 研究结论与展望}

\hangindent=2em \hangafter=1
6.1 研究结论

\hangindent=2em \hangafter=1
6.2 研究不足

\hangindent=2em \hangafter=1
6.3 未来展望

\noindent 参考文献

\noindent 附录

\vfill
\begin{flushright} （本表可附页） \end{flushright}
\end{UnifiedBox}

% ===================== 第 12-14 页：四、工作的重点与难点，拟采取的解决方案 =====================
\vspace{0.6cm}
\noindent {\zihao{4} \bfseries 四、工作的重点与难点，拟采取的解决方案}

\newpage

\begin{UnifiedBox}
\zihao{-4}
\setstretch{1.6}
\setlength{\parindent}{2em}

\noindent \textbf{（一）工作重点}

\vspace{0.2cm}
\noindent \textbf{1. 华润数科内外部战略环境的精准定性分析}

内外部环境分析是战略设计的前提，其精准性直接决定战略定位与路径的适配性，需重点突破三个维度的分析深度：一是政策环境的``机遇-约束''平衡分析，需梳理近5年国家及地方层面数字经济、国企数字化转型相关政策，区分``鼓励性政策''与``约束性政策''，明确政策对业务布局的具体影响；二是市场环境的``需求-竞争''匹配分析，需聚焦政企数字化转型细分领域，识别各领域的核心需求痛点与竞争格局，为华润数科找到``需求未被满足+竞争压力较小''的市场空间；三是内部资源能力的``优势-短板''客观梳理，需结合华润集团产业特性，明确可复用资源的``可及性''与``价值性''，同时正视市场化激励机制、核心技术自主可控性等短板，避免分析主观偏差。

\vspace{0.3cm}
\noindent \textbf{2. 集团协同型战略路径的适配性设计}

华润数科的核心竞争力源于集团产业资源，如何将资源优势转化为业务竞争力，是战略设计的核心重点。需重点设计两大协同路径：一是``集团内部数字化服务-外部行业拓展''的协同路径，明确内部服务与外部业务的资源复用机制，避免内部服务与外部业务形成资源争夺；二是``技术-场景-生态''的协同路径，结合工业互联网、数据安全等核心技术方向，匹配集团零售、能源等场景需求，同时筛选生态伙伴弥补技术短板，形成``技术支撑场景、场景吸引伙伴、伙伴完善生态''的闭环。此外，需分阶段设定战略目标，确保短期目标可落地、中期目标可实现、长期目标可支撑，避免目标与资源能力脱节。

\vspace{0.3cm}
\noindent \textbf{3. 国有科技企业特性适配的保障体系构建}

华润数科作为国有集团旗下科技企业，需平衡``国有资产监管要求''与``市场化经营灵活性''，这是保障战略落地的关键重点。需重点构建三大保障机制：一是合规保障机制，针对数据安全、国资增值等核心要求，制定覆盖业务全流程的合规审查流程，确保业务开展不触碰监管红线；二是市场化激励机制，参考科创板国企案例，设计差异化激励方案，解决国有科技企业``激励不足''问题；三是动态评估机制，基于平衡计分卡思路，设定贴合国有科技企业特性的评估指标，定期跟踪战略落地效果，及时调整执行策略，避免战略僵化。

\vspace{0.5cm}
\noindent \textbf{（二）工作难点及拟采取的解决方案}

\vspace{0.2cm}
\noindent \textbf{1. 难点一：集团内部资源协同机制的信息不对称问题}

\vspace{0.2cm}
\noindent \textbf{（1）难点表现}

华润集团业务覆盖零售、能源、医药等多个领域，各板块资源的管理权限、调用流程、收益分配规则分散在不同部门，公开信息有限，研究过程中易出现``资源可及性不明确''``协同成本难评估''等信息不对称问题，导致集团资源协同路径设计缺乏实操性。

\vspace{0.2cm}
\noindent \textbf{（2）拟采取的解决方案}

\noindent \textbf{多渠道信息补充：}通过华润集团官网下载年度报告、社会责任报告，梳理各板块业务布局与数字化需求；查阅集团内部发布的《资源共享管理办法》《数字化转型规划》等公开制度文件，明确资源调用的基本流程；同时，通过行业媒体报道、合作伙伴访谈，侧面了解集团资源协同的实际案例，补充公开信息缺口。

\noindent \textbf{专家咨询验证：}邀请熟悉华润集团运营模式的专家，对资源协同路径设计方案进行评审，判断资源调用流程的可行性、收益分配规则的合理性，避免因信息不对称导致方案脱离实际。

\noindent \textbf{分场景假设推演：}针对信息难以获取的资源，基于国有集团管理通用逻辑，设计``乐观-中性-悲观''三种场景假设，并分别制定适配的协同路径，确保方案在不同场景下均具备一定可行性。

\vspace{0.3cm}
\noindent \textbf{2. 难点二：国有与市场化特性平衡的激励机制设计边界模糊}

\vspace{0.2cm}
\noindent \textbf{（1）难点表现}

国有科技企业激励机制设计需同时满足``国资委监管要求''与``市场化激励效果''，但当前监管政策对``创新容错范围''``股权激励比例''等边界界定较为原则性，缺乏具体量化标准，导致激励机制设计易出现``过于保守''或``突破监管''的两难问题。

\vspace{0.2cm}
\noindent \textbf{（2）拟采取的解决方案}

\noindent \textbf{政策文本深度解读：}系统梳理《关于进一步完善国有企业中长期激励机制的通知》《国有企业科技人才薪酬分配指引》等监管文件，提取``激励对象范围''``激励额度上限''``容错认定标准''等核心条款，明确政策红线与弹性空间，确保机制设计不触碰监管底线。

\noindent \textbf{标杆案例对标分析：}选取国有科技企业市场化激励的典型案例，分析其激励对象选择、激励方式设计、合规审查流程，总结``监管要求-市场化需求''的平衡经验，为华润数科激励机制设计提供参考。

\noindent \textbf{分阶段试点推进：}设计``试点-推广''的渐进式方案，先在华润数科创新业务板块试点项目跟投机制，明确试点范围、跟投比例、容错条件，试点期结束后总结经验，优化机制设计，再逐步推广至其他业务板块，降低合规风险与实施阻力。

\vspace{0.3cm}
\noindent \textbf{3. 难点三：数字服务市场竞争格局的动态变化应对}

\vspace{0.2cm}
\noindent \textbf{（1）难点表现}

数字服务行业技术迭代快、市场需求变、竞争对手策略调整频繁，研究过程中基于当前格局设计的战略路径，可能因市场动态变化导致适配性下降，难以支撑长期发展。

\vspace{0.2cm}
\noindent \textbf{（2）拟采取的解决方案}

\noindent \textbf{动态竞争格局跟踪：}建立``月度行业信息汇总''机制，通过艾瑞咨询、IDC等行业数据库，跟踪数字服务市场的技术趋势、需求变化、竞争对手动态，及时更新市场竞争格局分析结论，为战略调整提供依据。重点关注人工智能、工业互联网、数据安全等核心技术领域的专利布局、融资动态与产品迭代信息，形成``技术趋势-市场需求-竞争策略''三维跟踪体系，确保战略研究能够及时反映市场最新变化。

\noindent \textbf{弹性战略路径设计：}在核心战略路径中预留``调整接口''，例如在技术创新路径中，明确``若某核心技术被替代，可快速切换至替代技术路线''；在市场拓展路径中，设定``若某细分领域竞争加剧，可优先聚焦集团关联行业客户''，避免战略路径僵化。同时，建立``战略路径-触发条件-替代方案''的映射表，将弹性设计制度化、文档化，确保在市场环境发生变化时能够快速响应、有序调整。

\noindent \textbf{多场景战略验证：}基于当前、1年后、3年后三种市场场景假设，分别验证战略路径的适配性，筛选出在多场景下均具备竞争力的核心路径，同时针对单一场景适配的路径，制定``场景触发-路径启动''的应对规则，提升战略对市场变化的应对能力。通过构建``基准场景-乐观场景-悲观场景''的分析框架，量化评估不同市场条件下战略目标的实现概率与资源需求，增强战略方案的稳健性与可操作性。

\vspace{8cm}

\begin{flushright} （本表可附页） \end{flushright}
\end{UnifiedBox}

% ===================== 第 15-16 页：五、论文工作量及进度 =====================
\newpage

\vspace{0.6cm}
\noindent {\zihao{4} \bfseries 五、论文工作量及进度}

\vspace{0.3cm}

% 第15-16页：工作量和进度规划都在同一个框内
\begin{UnifiedBox}[height=25.4cm]
\zihao{-4}
\setstretch{1.6}
\setlength{\parindent}{2em}

\noindent \textbf{（一）论文工作量}

本论文研究需完成``资料收集-分析研究-方案设计-报告撰写''全流程工作，总工作量约需180-200个标准工时，具体拆分如下：

\vspace{0.2cm}
\noindent \textbf{1. 资料收集与整理（40-45工时）}

\noindent \textbf{政策与行业资料收集：}需梳理2019-2024年国家及地方数字经济、国企数字化转型相关政策文件30+份，下载IDC、艾瑞咨询等权威机构数字服务行业报告15-20份，整理过程需对政策条款进行分类标注、对行业数据进行关键信息提取，预计耗时15-18工时。

\noindent \textbf{企业与案例资料收集：}收集华润集团及华润数科年度报告、数字化转型规划等公开资料10-12份，整理国有科技企业、产业数字服务厂商标杆案例8-10个，需对案例的战略路径、实施效果进行结构化梳理，预计耗时18-20工时。

\noindent \textbf{理论文献收集：}检索公司发展战略、国有科技企业转型相关中英文文献50-60篇，筛选核心文献25-30篇进行深度阅读，撰写文献综述提纲，预计耗时7-9工时。

\vspace{0.2cm}
\noindent \textbf{2. 内外部环境分析（35-40工时）}

\noindent \textbf{外部环境分析：}运用PEST分析法解读政策文件、行业报告，形成政策``机遇-约束''清单、市场``需求-竞争''分析报告，需多次调整分析维度以确保精准性，预计耗时15-18工时。

\noindent \textbf{内部资源能力分析：}基于华润数科公开资料，梳理资源与能力现状，总结优势与短板，形成内部分析报告，期间需反复验证信息真实性，预计耗时12-14工时。

\noindent \textbf{SWOT匹配分析：}结合内外部分析结果，构建SWOT矩阵，明确各象限策略方向，预计耗时8-10工时。

\vspace{0.2cm}
\noindent \textbf{3. 战略与保障体系设计（45-50工时）}

\noindent \textbf{战略定位与目标设计：}参考标杆案例，结合华润数科特性确定双重战略定位，分阶段设定目标，需多次研讨调整目标量化指标，预计耗时12-15工时。

\noindent \textbf{核心战略路径规划：}设计集团资源协同、技术创新、生态构建三大路径，需细化资源调用流程、技术研发步骤等实操内容，预计耗时18-20工时。

\noindent \textbf{保障体系构建：}制定合规、激励、评估三大保障机制，需查阅监管文件确认合规边界、参考案例设计激励方案，预计耗时15-17工时。

\vspace{0.2cm}
\noindent \textbf{4. 论文撰写与修改（60-65工时）}

\noindent \textbf{初稿撰写：}按照论文框架完成6章内容撰写，总字数预计6-8万字，需确保理论分析与实践设计逻辑连贯，预计耗时35-40工时。

\noindent \textbf{修改完善：}针对初稿进行内容补充、逻辑优化，同时调整格式规范，预计耗时15-18工时。

\noindent \textbf{定稿审核：}邀请导师对终稿进行评审，根据反馈意见进行最后修改，确保论文质量，预计耗时10-12工时。

\vspace{0.5cm}
\noindent \textbf{（二）论文进度规划}

\vspace{0.2cm}
本论文研究周期预计10个月，具体进度安排如下：

\vspace{0.3cm}
\noindent
\renewcommand{\arraystretch}{1.8}
\begin{tabular}{|>{\centering\arraybackslash}p{2.5cm}|>{\arraybackslash}p{5.5cm}|>{\arraybackslash}p{5.5cm}|}
\hline
\bfseries 时间节点 & \bfseries 核心任务 & \bfseries 交付成果 \\
\hline
第1-2个月 & 明确研究范围，完成政策、行业、企业、案例、理论文献收集与整理，撰写文献综述 & 文献综述初稿、资料汇编手册 \\
\hline
第3-4个月 & 开展外部环境（PEST、五力模型）与内部资源能力分析，完成SWOT匹配分析 & 内外部环境分析报告、SWOT矩阵图 \\
\hline
第5-6个月 & 确定战略定位与分阶段目标，设计集团协同型战略路径与三大保障机制 & 战略设计方案、保障机制细则 \\
\hline
第7-8个月 & 按照论文框架撰写初稿，完成章节内容整合与逻辑梳理 & 论文初稿 \\
\hline
第9个月 & 邀请导师、专家评审初稿，根据反馈进行内容补充、逻辑优化与格式调整 & 论文修改稿 \\
\hline
第10个月 & 针对修改稿进行最终审核，解决遗留问题，完成定稿与格式规范 & 论文终稿、答辩PPT \\
\hline
\end{tabular}

\vspace*{\fill}
\begin{flushright} （本表可附页） \end{flushright}
\end{UnifiedBox}

% ===================== 第 17-18 页：六、论文预期成果及创新点 =====================
\newpage

\vspace{0.6cm}
\noindent {\zihao{4} \bfseries 六、论文预期成果及创新点}

\vspace{0.3cm}

\begin{UnifiedBox}
\zihao{-4}
\setstretch{1.6}
\setlength{\parindent}{2em}


\vspace{0.2cm}
\noindent \textbf{1. 理论成果}

\noindent \textbf{国有集团数字子公司战略分析框架：}基于动态能力理论、生态系统理论，结合国有科技企业特性，构建``外部环境-内部资源-战略定位-路径设计-保障体系''五维定性分析框架，填补现有研究对``国有集团背景+数字服务业务''企业战略设计的理论空白，为同类企业的战略研究提供理论参考。

\noindent \textbf{集团资源协同定性分析方法：}针对国有集团内部资源信息不对称问题，提出``多渠道信息补充+专家咨询验证+分场景假设推演''的定性分析方法，明确资源协同路径设计的核心逻辑，丰富国有集团资源管理的理论方法体系。

\vspace{0.2cm}
\noindent \textbf{2. 实践成果}

\noindent \textbf{华润数科集团协同型发展战略方案：}明确``集团数字化转型核心载体+行业数字服务生态构建者''的双重定位，分短期（1-2年）、中期（3-5年）、长期（5年以上）设定可落地的战略目标，设计集团资源协同、技术创新、生态构建三大核心路径，为华润数科制定具体业务规划提供直接指导。

\noindent \textbf{国有科技企业战略实施保障机制细则：}制定覆盖合规、激励、评估的三大保障机制细则，包括数据安全合规审查流程、创新业务板块项目跟投方案、国有科技企业特色评估指标体系，可直接应用于华润数科战略落地，同时为其他国有科技企业提供可复用的保障机制模板。

\vspace{0.2cm}
\noindent \textbf{3. 应用价值}

\noindent \textbf{对华润数科的价值：}帮助企业清晰认知内外部环境，明确战略发展方向与实操路径，解决集团资源协同、市场化激励等核心问题，支撑企业实现``内部数字化服务覆盖+外部行业市场拓展''的双轮驱动发展。

\noindent \textbf{对国有科技企业的价值：}为国有集团旗下数字科技企业提供战略设计与实施的参考案例，尤其在``国有与市场化平衡''''集团资源复用''等关键问题上，提供可借鉴的思路与方法。

\noindent \textbf{对数字服务行业的价值：}梳理政企数字化转型细分领域的需求痛点与竞争格局，总结国有科技企业在数字服务市场的差异化竞争优势，为行业参与者理解市场竞争态势提供参考。

\vspace{0.5cm}
\noindent \textbf{（二）论文创新点}

\vspace{0.2cm}
\noindent \textbf{1. 研究视角创新：聚焦国有集团数字子公司的``双重属性''平衡}

现有研究多单独关注``国有科技企业的监管合规''或``数字服务企业的市场化发展''，较少针对``国有集团旗下数字子公司''的``集团资源依赖+国有监管约束+市场化竞争压力''三重特性展开研究。本论文首次系统分析华润数科的``双重属性''，设计``政策响应-资源协同-市场竞争''三位一体的战略框架。

\vspace{0.2cm}
\noindent \textbf{2. 战略设计创新：构建``集团资源-场景需求-生态伙伴''协同路径}

突破传统数字服务企业``技术导向''的战略设计逻辑，基于华润集团零售、能源等产业场景积淀，设计``集团资源复用-场景需求匹配-生态伙伴补位''的协同路径，形成``从集团内部场景到外部行业市场''的业务拓展模式，既发挥国有集团资源优势，又解决数字服务企业``场景落地难''的核心问题。

\vspace{0.2cm}
\noindent \textbf{3. 保障机制创新：提出国有科技企业``渐进式市场化激励''方案}

针对国有科技企业激励机制``边界模糊''问题，参考科创板国企案例，创新设计``分阶段试点-动态调整''的渐进式激励方案：第一阶段在创新业务板块试点项目跟投，明确小范围、低比例的激励规则；第二阶段根据试点效果，逐步扩大激励范围、优化激励方式，同时严格对接国资委监管要求，解决``激励不足''与``合规风险''的两难问题，为国有科技企业激励机制设计提供创新思路。

\vspace*{\fill}

\begin{flushright} （本表可附页） \end{flushright}
\end{UnifiedBox}

% ===================== 第 19-22 页：七、主要参考文献 =====================
\newpage

\vspace{0.6cm}
\noindent {\zihao{4} \bfseries 七、主要参考文献}

\vspace{0.3cm}

\begin{UnifiedBox}[height=25.4cm]
\zihao{5}
\setstretch{1.35}
\setlength{\parindent}{0em}
\setlength{\leftmargini}{1.5em}

\noindent （要求35篇，文献阅读需达到32学时）

\vspace{0.3cm}

\begin{enumerate}[label={[}\arabic*{]}, leftmargin=2em, itemsep=0.15em, parsep=0pt]

\item Lin W. The Sustainable Development Strategy Research of Small and Medium-Sized Enterprises in the Economic Transition Period[J]. Industrial Engineering and Innovation Management, 2024, 7(3): DOI:10.23977/IEIM.2024.070310.

\item 李亚祥. 基于SWOT分析法的铁路物资公司发展战略研究[J]. 铁路采购与物流, 2024, 19(04): 24-26.

\item 祝志华. 汽车零部件公司发展战略研究——以T公司为例[J]. 经济师, 2024, (04): 287-288.

\item Zhang D W. Commercial Company Development Strategy Research and Exploration[J]. The Journal of Education Insights, 2024, 2(1).

\item 谢飞. YS公司发展战略研究[J]. 经济师, 2023, (03): 275-276.

\item 汪恭二, 占寿罡, 万超. 基于SWOT分析法的T公司发展战略研究[J]. 铜业工程, 2023, (01): 196-200.

\item 吴伟泽. 基于SWOT分析的G股份有限公司发展战略研究[J]. 北方经贸, 2022, (03): 135-137.

\item 龙敏艳. 公司发展战略分析——以华为为例[J]. 老字号品牌营销, 2022, (05): 21-23.

\item 张如, 赵心坦. 基于SWOT分析的J集团公司发展战略研究[J]. 中国乡镇企业会计, 2021, (09): 118-120.

\item 杨华江. 公司战略和战略风险管理理论的演进研究[J]. 科学决策, 2021, (07): 124-135.

\item Vochozka H. Scientific and Technological Innovation and Enterprise Sustainable Development Strategy Research[J]. Advance in Sustainability, 2021, 1(1).

\item 袁琪. 大连机场货运公司发展战略研究——基于五力模型分析[J]. 空运商务, 2021, (06): 50-55.

\item 蔡洋. SWOT视角下蓝月亮实业有限公司发展战略研究[J]. 现代营销(下旬刊), 2021, (12): 130-131.

\item 王广昀, 王凤兰, 赵波, 等. 大庆油田公司勘探开发形势与发展战略[J]. 中国石油勘探, 2021, 26(01): 55-73.

\item 孙虹. 基于IFE矩阵和EFE矩阵的X公司的发展战略研究[J]. 中小企业管理与科技(中旬刊), 2020, (10): 78-80.

\item 叶澜杰, 张象麟. ``两票制''背景下医药公司发展战略研究[J]. 现代商业, 2020, (23): 35-36.

\item 权小锋, 醋卫华, 徐星美. 高管从军经历与公司盈余管理：军民融合发展战略的新考察[J]. 财贸经济, 2019, 40(01): 98-113.

\item 陈燕丽, 王磊, 姜明栋, 等. 东北三省制造业上市公司企业绩效及影响因素研究——基于DEA-Malmquist-Tobit模型[J]. 工业技术经济, 2018, 37(11): 51-57.

\item Clayton A P. Innovation in local economic development strategy: The Wireless Research Center of North Carolina[J]. Geography Compass, 2018, 12(6): n/a-n/a.

\item 金之钧, 蔡勋育, 刘金连, 等. 中国石油化工股份有限公司近期勘探进展与资源发展战略[J]. 中国石油勘探, 2018, 23(01): 14-25.

\item Murtini S, Kuspriyanto, Kurniawati A. Mangrove area development strategy wonorejo as ecotourism in surabaya[J]. Journal of Physics: Conference Series, 2018, 953(1): 012174-012174.

\item Dong X. Development strategy research of low-carbon tourist city[J]. IOP Conference Series: Earth and Environmental Science, 2017, 61(1): 012158-012158.

\item 陈凯麟, 蒋伏心. ``供给侧''改革背景下的证券公司发展战略研究——基于互联网金融模式下的GL证券公司实证研究[J]. 经济体制改革, 2017, (02): 100-106.

\item 巩娜, 郭翠菱. 民营高科技企业战略转型与公司持续发展[J]. 经济体制改革, 2017, (01): 116-122.

\item Bing L, Xianghua X, Hua X, et al. Zhengdong New District Smart Government Application System Construction and Its Development Strategy Research[J]. MATEC Web of Conferences, 2017, 10005069.

\item Shuai W, Hongxiang C. Research on the development strategy of Pingdingshan iron and steel industry cluster with Wugang Company as the core[J]. MATEC Web of Conferences, 2017, 10005066-05066.

\item 韩东. 基于发展战略视角的企业财务报表分析——以美诺华药业股份有限公司为例[J]. 财会通讯, 2016, (35): 75-79.

\item Lysogene; Lysogene Expands Senior Leadership in the USA to Support Its International Development Strategy[J]. Biotech Week, 2016.

\item Liu X. China new pattern urbanization process medium and small towns sports culture development strategy research[J]. BioTechnology: An Indian Journal, 2014, 10(8).

\item Yao W. Analytic hierarchy process-based recreational sports events development strategy research[J]. BioTechnology: An Indian Journal, 2014, 10(6).

\item Zhang B. City spontaneous sports organizations development strategy research in the perspective of public satisfaction index[J]. BioTechnology: An Indian Journal, 2014, 10(9).

\item Zhang P, Sun J. Chinese new pattern urbanization process medium and small towns sports industrial development strategy research[J]. BioTechnology: An Indian Journal, 2014, 10(11).

\item Song F. Goal programming-based China sports tourism group preference and development strategy research[J]. BioTechnology: An Indian Journal, 2014, 10(15).

\item 李长娥, 谢永珍. 区域经济发展水平、女性董事对公司技术创新战略的影响[J]. 经济社会体制比较, 2016, (04): 120-131.

\item YanYu. People love and sports tourism industry development strategy research[J]. Journal of Residuals Science \& Technology, 2016, 13(8).

\item 刘建刚, 钱玺娇. ``互联网+''战略下企业技术创新与商业模式创新协同发展路径研究——以小米科技有限责任公司为案例[J]. 科技进步与对策, 2016, 33(01): 88-94.

\item 李玲. 我国金融不良资产的发展趋势、监管政策与处置机制——兼论大型资产管理公司的战略取向[J]. 新金融, 2015, (11): 38-44.

\item Xiangyang Z, Shaohua Y. Deaf-Mute Students Physical Health Development Strategy Research in the Perspective of Education Fair Equalization[J]. The Open Cybernetics \& Systemics Journal, 2015, 9(1): 1640-1645.

\item Huo J, Han R. Sustainable Development Strategy Research of Historical Village Environment - Take Huaiyin Village of Zhenjiang as Example[J]. Applied Mechanics and Materials, 2014, 3547(675-677): 1246-1252.

\item Li P F. Sanda Market Operation and Industrialization Development Strategy Research[J]. Advanced Materials Research, 2014, 3255(971-973): 2406-2409. DOI:10.4028/www.scientific.net/AMR.971-973.2406.

\item 崔和瑞, 葛静. 基于SWOT-PEST分析范式的新疆电力公司发展战略研究[J]. 电力学报, 2013, 28(03): 203-210.

\item Lee, Wan-Bok, Park, et al. Festival Program Development Strategy Research for the Activation of the Local Festivals: Focusing on the Osan Doksanseong Cultural Festival[J]. Humanities Contens, 2013, (28): 219-241.

\item 郑晓明, 丁玲, 欧阳桃花. 双元能力促进企业服务敏捷性——海底捞公司发展历程案例研究[J]. 管理世界, 2012, (02): 131-147+188.

\item Junjie C, Ming L, Shuguo L. Development Strategy Research of Modern Eco-Agriculture on the basis of constructing the Rural Circular Economy-For the Example of Shandong Province[J]. Energy Procedia, 2011, 5(C): 2504-2508.

\item 姚铮, 汤彦峰. 商业银行引进境外战略投资者是否提升了公司价值——基于新桥投资收购深发展的案例分析[J]. 管理世界, 2009, (S1): 94-102+133.

\item 芮明杰, 詹文静, 陈杰. 跨区域发展战略对房地产企业绩效的影响——基于房地产上市公司的实证研究[J]. 中国工业经济, 2008, (08): 56-64.

\item 李卓, 刘杨, 陈永清. 发展中国家跨国公司的国际化战略选择：针对中国企业实施``走出去''战略的模型分析[J]. 世界经济, 2006, (11): 11-22+95.

\item 吴文武. 跨国公司与经济发展——兼论中国的跨国公司战略[J]. 经济研究, 2003, (06): 38-44+94.

\item 赵景华. 跨国公司在华子公司成长与发展战略的实证研究[J]. 管理世界, 2002, (10): 93-101.

\item 赵景华. 跨国公司在华子公司成长与发展的战略角色及演变趋势[J]. 中国工业经济, 2001, (12): 61-66.

\end{enumerate}

\vspace*{\fill}
\begin{flushright} （本表可附页） \end{flushright}
\end{UnifiedBox}
\vfill %
\end{document}
